\section{Introduction} 

The process of choosing a career is undeniably one of the most critical and defining decisions a student will make in their lifetime. This single choice not only dictates their higher education trajectory but also establishes the foundation for their socioeconomic mobility, personal fulfillment, and contribution to national development. However, despite the gravity of this decision, many students make this choice based on highly constrained and often skewed information. In developing nations particularly, peer pressure, societal expectations, and parental demands frequently overshadow a clear, objective understanding of the student's own academic strengths, weaknesses, and intrinsic aptitudes. This misalignment is particularly detrimental when attempting to guide students toward the health sector, an arena that demands absolute precision, resilience, and specific, deep foundational knowledge in core science subjects such as Biology, Chemistry, Mathematics, and Physics.

To address this systemic challenge, researchers and educators have increasingly turned toward intelligent guidance systems. These systems leverage the vast amounts of historical data generated by educational institutions to analyze a student's past academic performance and subsequently recommend highly suitable educational pathways and careers. However, traditional systems deployed in educational data mining have often treated this complex, multifaceted challenge as a simple, mutually exclusive classification problem. They provide a singular, absolute recommendation (e.g., "You should become a Doctor") without exploring alternative viable paths and, crucially, without mathematically or logically explaining \textit{why} the recommendation was made. 

This thesis introduces a comprehensive conceptual framework designed to revolutionize this approach. The framework models the student journey not as a single point of decision, but as a continuous, intelligent, and interconnected pathway:
\[
\boxed{\text{Academic Performance (Input)} \rightarrow \text{Educational Pathway} \rightarrow \text{Career Pathway (Output)}}
\]
Specifically, this study develops an Explainable Multi-Label \parencite{tsoumakas2007multi} Machine Learning Model designed explicitly to assist students in Rwanda transitioning from their Ordinary Level (Senior 3 or S3) to Advanced Level (S6) and, ultimately, into highly specialized health sector careers. The system proposed herein integrates multiple layers of artificial intelligence, including prediction, multi-label \parencite{tsoumakas2007multi} recommendation, and explicit mathematical explanation, moving far beyond simple black-box algorithmic models to provide actionable, ranked, and logically transparent guidance that students and educators can actively trust.

\subsection{Background of the Study}

In the Rwandan education system, the academic journey is sharply bifurcated at critical assessment junctures. Students complete their Ordinary Level (S3) education by taking rigorous national examinations administered by the National Examination and School Inspection Authority (NESA). The results from these national exams, particularly in the core STEM subjects (Science, Technology, Engineering, and Mathematics), serve as the strongest, most objective indicators of a student's cognitive aptitude and academic preparedness. The transition from S3 to Advanced Level (S6) is not merely a change in grade; it is a definitive filtering mechanism where students are channeled into highly specialized subject combinations (such as PCB: Physics-Chemistry-Biology, or MCB: Mathematics-Chemistry-Biology). These combinations strictly dictate their eligibility for future university degrees and directly constrain their options in the labor market. If a student with a high aptitude for healthcare mistakenly chooses a combination lacking Biology or Chemistry due to poor guidance, their path to becoming a medical professional is essentially severed.

The health sector represents a critical, high-priority area for Rwanda's national development. Under the framework of Rwanda Vision 2050 and the National Strategy for Transformation (NST1), the government has explicitly prioritized the development of a knowledge-based economy and the establishment of world-class healthcare infrastructure. To achieve these ambitious macroeconomic and social goals, the nation requires a steady, high-quality, and predictable pipeline of capable professionals, including specialized physicians, advanced practice nurses, biomedical engineers, pharmacists, and public health administrators. Consequently, guiding the right students, those possessing the precise mathematical, analytical, and scientific aptitudes required to survive rigorous medical training, into the appropriate educational pathways is not just an individual necessity, but a macroeconomic imperative that directly impacts national mortality rates and healthcare efficiency.

Historically, this critical guidance has been conducted manually by educators, school administrators, and career counselors. However, human-led guidance is inherently limited in scope and accuracy. It is highly susceptible to cognitive biases, constrained by the very limited time available per student, and mathematically incapable of processing the massive, multi-dimensional historical datasets required to recognize subtle, complex patterns of long-term academic success and failure. For instance, a human counselor cannot instantly calculate the precise probability of a student succeeding in Medical School given an 82\% in Chemistry, a 74\% in Physics, and a 91\% in Biology compared against the historical outcomes of ten thousand similar students. 

Machine learning (ML) offers a profound, scalable solution by autonomously identifying these complex, non-linear patterns in large historical datasets. Yet, despite the immense promise of Artificial Intelligence, typical ML approaches utilized in educational data mining and career prediction harbor significant methodological and theoretical limitations that prevent widespread adoption:

\begin{enumerate}
    \item \textbf{Single-Label Constraint:} Traditional models almost exclusively predict only a single career outcome. They frame the complex human capability matrix as a mutually exclusive choice (e.g., classifying a student strictly as either a Doctor OR a Nurse). This ignores the fundamental reality of human aptitude, that a student's unique skills might make them highly suitable for multiple related health sector careers simultaneously.
    \item \textbf{The Black-Box Dilemma:} Advanced models, particularly deep neural networks and complex ensembles, operate as opaque "black boxes." They provide a final prediction without any transparent mathematical or logical reasoning. In a high-stakes domain like educational tracking, providing a prediction without an explanation destroys user trust. If a student is told by an algorithm that they are statistically unlikely to succeed in Medicine, they must mathematically understand exactly which specific academic deficiencies led to that conclusion in order to accept it or attempt to improve.
    \item \textbf{Static Thresholding:} Many legacy guidance systems rely on hard-coded static thresholds (e.g., "Must have $>80\%$ in Biology to study Medicine"), failing to account for the dynamic interactions and covariance between different STEM subjects that a machine learning model inherently captures.
\end{enumerate}

\subsection{Problem Statement}

Despite the modern availability of vast, rich amounts of educational data (such as the comprehensive S3 national exam datasets collected from 2017 to 2025 across Rwanda), there is a severe lack of advanced, multi-label \parencite{tsoumakas2007multi}, and explainable machine learning models utilized for career path guidance. 

Existing predictive models generally treat career guidance as a mutually exclusive choice. By framing the algorithm as a standard multi-class problem, they attempt to classify a student strictly as a future "Doctor" OR a "Nurse" OR a "Pharmacist." This approach is fundamentally flawed because it forces the algorithm to suppress secondary, highly viable options, thereby failing to output ranked probabilities across multiple pathways that the student might excel in.

Furthermore, these traditional systems completely lack the critical intelligence layer of explainability. When a student is strongly recommended to pursue Pharmacy, current predictive systems cannot mathematically, visually, or logically explain to the student or their parents which specific exam scores contributed most heavily to that specific recommendation. They cannot answer the "Why?" 

This study rigorously addresses these empirical and theoretical gaps by implementing a multi-label \parencite{tsoumakas2007multi} classification approach structurally combined with cutting-edge Explainable AI (XAI) techniques. The result is a system that not only predicts multiple career pathways simultaneously, acknowledging the multifaceted nature of student capability, but also explicitly explains the mathematical reasoning driving those recommendations.

\subsection{Objective of the Study}

\textbf{General Objective:} \\
To research, develop, and evaluate an explainable, multi-label \parencite{tsoumakas2007multi} machine learning model capable of accurately predicting educational and career pathways in the health sector based strictly on Rwandan S3 academic performance data.

\textbf{Specific Objectives:}
\begin{itemize}
    \item To fundamentally restructure career prediction from a single-label to a multi-label \parencite{tsoumakas2007multi} classification paradigm by implementing the Classifier Chains \parencite{read2011classifier} mechanism.
    \item To train, tune, and comparatively evaluate multiple base machine learning algorithms, specifically Extreme Gradient Boosting (XGBoost), Random Forest \parencite{breiman2001random}, Decision Trees, and Multi-Layer Perceptrons (MLP).
    \item To generate ranked, probabilistic outputs for multiple career outcomes rather than a single, absolute, and limiting prediction.
    \item To integrate a robust explainability layer utilizing feature importance mapping to explicitly identify and quantify which academic subjects drive specific career recommendations.
    \item To explore and formulate the fundamental mathematics of counterfactual \parencite{wachter2017counterfactual} explanations, providing a theoretical mechanism to show students exactly what academic score changes would alter their career recommendations.
\end{itemize}

\subsection{Research Questions}

To achieve the aforementioned objectives, this study seeks to systematically answer the following research questions:
\begin{enumerate}
    \item How effectively can a multi-label \parencite{tsoumakas2007multi} classification approach, specifically Classifier Chains \parencite{read2011classifier}, model the non-mutually exclusive, highly correlated nature of health sector career pathways compared to traditional single-label methods?
    \item How do advanced tree-based ensemble models (XGBoost, Random Forest) compare to foundational models (Decision Tree) and deep learning models (MLP) in terms of Subset Accuracy and Hamming Loss when predicting multi-label \parencite{tsoumakas2007multi} career paths based on exam scores?
    \item In what ways can Explainable AI (XAI) techniques mathematically clarify the complex, non-linear relationships between specific S3 subject scores and the resulting multi-label \parencite{tsoumakas2007multi} career recommendations?
\end{enumerate}

\subsection{Hypotheses}

Based on the theoretical framework and existing literature in the broader field of machine learning, this research proposes the following testable hypotheses:

\begin{itemize}
    \item \textbf{H1 (Architectural Superiority):} Multi-label models combined with the Classifier Chains \parencite{read2011classifier} mechanism will yield significantly higher predictive subset accuracy and lower Hamming Loss for career recommendations than completely independent, unchained binary classifiers (Binary Relevance), because they capture the inherent label correlations between related health careers.
    \item \textbf{H2 (Algorithmic Dominance):} Ensemble tree-based methods, specifically XGBoost \parencite{chen2016xgboost} and Random Forest \parencite{breiman2001random}, will consistently outperform singular Decision Trees and basic Multi-Layer Perceptrons in capturing the specific complex, non-linear, and non-normalized relationships present in raw student grading data.
    \item \textbf{H3 (Interpretability Enhancement):} The structural integration of an explainability layer will successfully extract global feature importance, allowing for a clear, verifiable mathematical mapping between the input features (academic scores in Biology, Chemistry, Math, Physics) and the model's prediction output, thereby resolving the black-box dilemma.
\end{itemize}

\subsection{Significance of the Study}

The primary significance of this research lies in its paradigm shift away from basic, isolated algorithmic comparison and toward the creation of a comprehensive, human-centric, and highly intelligent educational framework. By explicitly prioritizing mathematical explainability alongside multi-label \parencite{tsoumakas2007multi} probability ranking, this study provides a highly practical, deployable tool that addresses the exact needs of various stakeholders in the Rwandan educational ecosystem.

\subsubsection{Impact on Students and Parents}
For students and parents, the system offers unprecedented transparent guidance, mitigating the anxiety and uncertainty that traditionally surround the S3 to S6 transition. If a student is recommended for a Bachelor of Pharmacy with a high probability of $P = 0.85$, the model's explainability component will mathematically clarify that this is largely driven by their exceptionally high scores in Chemistry and Mathematics, which successfully offset a slightly lower score in Biology. Conversely, if a student is rejected from a Medical track, the counterfactual \parencite{wachter2017counterfactual} engine provides actionable feedback (e.g., "An increase of 8\% in Physics would have changed this recommendation"). This transparency empowers students to make highly informed, logical decisions about their advanced educational pathways rather than relying on guesswork, societal pressure, or subjective human opinions.

\subsubsection{Impact on Educational Institutions and Counselors}
For educators, headteachers, and school counselors, the model acts as a highly advanced decision-support tool. It has the computational power to process thousands of historical data points instantly, effectively multiplying the counselor's capacity. By automating the objective, data-driven baseline recommendations, counselors are freed from manual transcript analysis. This allows them to focus their limited time entirely on the human, psychological, and emotional aspects of guidance, such as discussing a student's financial constraints, personal passions, or family expectations, while confidently relying on the algorithm for the empirical academic baseline.

\subsubsection{Impact on National Healthcare Infrastructure and Policy}
At a macroeconomic level, this model clearly demonstrates how historical national examination data, often left dormant in government databases after graduation, can be effectively leveraged as a form of predictive intelligence. By optimizing the pipeline of talent entering the critical health sector, this system ensures that university medical and nursing programs receive candidates who possess the precise, verified cognitive aptitudes required to succeed. This reduces university dropout rates, increases the quality of graduating healthcare professionals, and directly supports the long-term healthcare infrastructure goals outlined in Rwanda's National Strategy for Transformation (NST1).

\subsection{Scope of the Study}

To maintain a rigorous, mathematically sound, and feasible focus suitable for a Master's thesis, it is crucial to clearly define the boundaries of this research within the much broader conceptual vision of AI in education. This thesis explicitly divides the grand vision into current implementations and future doctoral-level extensions.

\textbf{Included in this Thesis (The Master's Scope):}
\begin{itemize}
    \item \textbf{Prediction \& Recommendation:} The complete implementation of multi-label \parencite{tsoumakas2007multi} classification to output mathematically ranked probabilities for various health careers based on cross-sectional historical data.
    \item \textbf{Explanation:} The implementation of feature importance metrics to explain the global behavior of the models.
    \item \textbf{Counterfactual Formulation:} Establishing the strict mathematical foundations and theoretical optimization formulas for basic counterfactual \parencite{wachter2017counterfactual} explanations (e.g., $x^* = \arg\min_{x'} d(x,x')$ subject to $f(x') \neq f(x)$).
    \item \textbf{Algorithms Evaluated:} XGBoost \parencite{chen2016xgboost}, Random Forest \parencite{breiman2001random}, Decision Tree, and MLP using the Classifier Chains \parencite{read2011classifier} methodology.
\end{itemize}

\textbf{Excluded from this Thesis (Future Research / The PhD Vision):}
\begin{itemize}
    \item \textbf{Continuous Evaluation (Dynamic Adaptation):} The proposed system will be trained on static historical data. It will not dynamically update its mathematical weights in real-time as new students are assessed ($R_t = f(X_t)$). Achieving this requires an active, longitudinal feedback loop outside the current scope.
    \item \textbf{Longitudinal Modeling:} Tracking the exact same cohort of students over many years through university to verify their actual final career placement against the S3 prediction.
    \item \textbf{Labour-Market Integration:} Incorporating external, fluctuating economic data, salary trends, or real-time hospital job market demand into the recommendation logic.
    \item \textbf{Advanced Career Graphing:} Building sophisticated Graph Neural Networks (GNN) to explicitly map the complex, multi-node subject-combination-career relationships as a mathematical graph.
\end{itemize}

By strictly defining these boundaries, this thesis establishes a highly ambitious yet technically achievable framework: a rigorous current implementation that serves as the concrete, proven foundation for expansive future extensions.

\subsection{Limitations and Assumptions of the Study}

While this research aims to provide a robust, highly accurate predictive model, it is fundamentally constrained by several empirical limitations and underlying mathematical assumptions that must be explicitly acknowledged:

\subsubsection{Assumptions}
\begin{itemize}
    \item \textbf{Cognitive Reflection in Grading:} A foundational assumption of this model is that a student's performance on the S3 national examinations in STEM subjects is a direct, objective, and accurate reflection of their cognitive aptitude for those subjects. It assumes the grading mechanism is perfectly standardized and devoid of systemic errors.
    \item \textbf{Stationarity of the Education System:} The model operates under the mathematical assumption of stationarity, meaning it assumes that the difficulty of the national exams, the grading rubrics, and the academic requirements for university health programs remain relatively constant between 2017 and 2025.
\end{itemize}

\subsubsection{Limitations}
\begin{itemize}
    \item \textbf{Exclusion of Socio-Economic Variables:} The current model relies exclusively on quantitative academic scores. It does not possess data regarding a student's socio-economic status, geographical location, psychological resilience, or financial capacity to afford medical school. These hidden variables ($Z$) introduce unmeasured variance into the real-world probability of a student actually achieving the recommended career.
    \item \textbf{Static Data Snapshot:} The model is trained on a static, cross-sectional snapshot of historical data. It cannot track an individual student's intellectual growth or decline between S3 (Ordinary Level) and S6 (Advanced Level).
    \item \textbf{Subject Scope Constraint:} To maintain high predictive accuracy specifically for the health sector, the model's feature space is strictly constrained to Biology, Chemistry, Mathematics, and Physics. It ignores humanities or language scores, which, while generally less critical for clinical medicine, still play a role in a student's overall academic profile.
\end{itemize}

\subsection{Operational Definition of Terms}
To ensure clarity throughout the technical discussions in the following chapters, key terms are defined operationally as follows:
\begin{itemize}
    \item \textbf{Multi-Label Classification:} A machine learning paradigm where each input instance (a student) can be assigned multiple target labels simultaneously (e.g., a student is mathematically classified as suited for both Medicine and Nursing). This contrasts with multi-class classification, where only one label can be chosen.
    \item \textbf{Explainable AI (XAI):} A suite of mathematical methods and algorithms designed to ensure that the internal mechanics, feature weights, and final results of an artificial intelligence solution can be thoroughly understood, audited, and trusted by human operators.
    \item \textbf{Classifier Chains (CC):} An advanced meta-algorithm for multi-label \parencite{tsoumakas2007multi} classification that links several binary classifiers in a sequential chain. Each classifier in the chain includes the predictions of all previous classifiers as additional input features, thereby allowing the model to capture deep statistical correlations between labels.
    \item \textbf{Counterfactual Explanation:} A specific form of local explanation that mathematically describes the absolute smallest theoretical change to the input data (e.g., increasing a Mathematics score by 4 points) that would result in the algorithm flipping its output to a different prediction.
    \item \textbf{Ensemble Learning:} A machine learning technique that strategically combines several base models (such as hundreds of Decision Trees) in order to produce one highly optimal, generalized predictive model that reduces variance (as in Random Forests) or bias (as in XGBoost).
\end{itemize}
